% ============================================================
% MANUAL DE USUARIO — Sistema de Gestión de Mantenimiento
% de Equipos (SGME)
% ============================================================
\documentclass[12pt, a4paper]{article}

% ── Paquetes ─────────────────────────────────────────────────
\usepackage[utf8]{inputenc}
\usepackage[spanish]{babel}
\usepackage[T1]{fontenc}
\usepackage{geometry}
\usepackage{graphicx}
\usepackage{xcolor}
\usepackage{listings}
\usepackage{booktabs}
\usepackage{tabularx}
\usepackage{hyperref}
\usepackage{float}
\usepackage{fancyhdr}
\usepackage{enumitem}
\usepackage{tcolorbox}
\usepackage{multicol}

% ── Configuración de página ──────────────────────────────────
\geometry{margin=2.5cm}

% ── Colores ──────────────────────────────────────────────────
\definecolor{primary}{HTML}{6C5CE7}
\definecolor{secondary}{HTML}{00CEC9}
\definecolor{darkbg}{HTML}{1A1A2E}
\definecolor{codebg}{HTML}{F5F5F5}
\definecolor{codeframe}{HTML}{DDDDDD}
\definecolor{accentgreen}{HTML}{00B894}
\definecolor{accentorange}{HTML}{F39C12}
\definecolor{accentred}{HTML}{E74C3C}

% ── Estilo de código ─────────────────────────────────────────
\lstset{
    backgroundcolor=\color{codebg},
    frame=single,
    rulecolor=\color{codeframe},
    basicstyle=\ttfamily\small,
    keywordstyle=\color{primary}\bfseries,
    commentstyle=\color{accentgreen}\itshape,
    stringstyle=\color{accentorange},
    breaklines=true,
    numbers=none,
    tabsize=4,
    showstringspaces=false,
    captionpos=b,
    xleftmargin=0.5em,
    framexleftmargin=0.5em,
    aboveskip=0.8em,
    belowskip=0.8em,
    literate={á}{{\'a}}1 {é}{{\'e}}1 {í}{{\'i}}1 {ó}{{\'o}}1 {ú}{{\'u}}1
             {ñ}{{\~n}}1 {Á}{{\'A}}1 {É}{{\'E}}1 {Í}{{\'I}}1 {Ó}{{\'O}}1
             {Ú}{{\'U}}1 {Ñ}{{\~N}}1,
}

\lstdefinestyle{console}{
    backgroundcolor=\color{darkbg},
    basicstyle=\ttfamily\small\color{white},
    frame=single,
    rulecolor=\color{darkbg},
    xleftmargin=0.5em,
    framexleftmargin=0.5em,
}

% ── Hipervínculos ────────────────────────────────────────────
\hypersetup{
    colorlinks=true,
    linkcolor=primary,
    urlcolor=secondary,
}

% ── Encabezados ──────────────────────────────────────────────
\pagestyle{fancy}
\fancyhf{}
\fancyhead[L]{\small\textcolor{gray}{SGME --- Manual de Usuario}}
\fancyhead[R]{\small\textcolor{gray}{\thepage}}
\fancyfoot[C]{\small\textcolor{gray}{Sistema de Gesti\'on de Mantenimiento de Equipos}}
\renewcommand{\headrulewidth}{0.4pt}

% ── Cajas ────────────────────────────────────────────────────
\newtcolorbox{infobox}[1][]{
    colback=primary!5, colframe=primary!60,
    fonttitle=\bfseries, title=#1,
    arc=3mm, boxrule=0.5mm
}

\newtcolorbox{tipbox}[1][]{
    colback=accentgreen!5, colframe=accentgreen!60,
    fonttitle=\bfseries, title=#1,
    arc=3mm, boxrule=0.5mm
}

\newtcolorbox{warnbox}[1][]{
    colback=accentorange!5, colframe=accentorange!60,
    fonttitle=\bfseries, title=#1,
    arc=3mm, boxrule=0.5mm
}

\newtcolorbox{errorbox}[1][]{
    colback=accentred!5, colframe=accentred!60,
    fonttitle=\bfseries, title=#1,
    arc=3mm, boxrule=0.5mm
}

% ══════════════════════════════════════════════════════════════
\begin{document}

% ── Portada ──────────────────────────────────────────────────
\begin{titlepage}
    \centering
    \vspace*{2cm}

    {\Huge\bfseries\textcolor{primary}{Manual de Usuario}\par}
    \vspace{0.5cm}
    {\Large\textcolor{secondary}{Sistema de Gesti\'on de Mantenimiento de Equipos}\par}
    \vspace{0.3cm}
    {\large SGME v1.0\par}

    \vspace{1.5cm}

    \rule{12cm}{2pt}

    \vspace{1.5cm}

    {\large Gu\'ia completa para el uso del sistema,\\
    incluyendo ejemplos pr\'acticos y soluci\'on de errores.\par}

    \vspace{2cm}

    {\large\textbf{Presentado por:}\par}
    \vspace{0.4cm}
    {\Large Emanuel Cardenas\par}
    \vspace{0.2cm}
    {\Large Samuel Manchola\par}
    \vspace{0.2cm}
    {\Large Sebastian Medina\par}

    \vfill

    \rule{12cm}{1pt}
    \vspace{0.5cm}
    {\large Mayo 2026\par}
\end{titlepage}

% ── Tabla de Contenidos ──────────────────────────────────────
\tableofcontents
\newpage

% ══════════════════════════════════════════════════════════════
% SECCIÓN 1: CÓMO INICIAR EL SISTEMA
% ══════════════════════════════════════════════════════════════
\section{C\'omo Iniciar el Sistema}

\subsection{Requisitos Previos}

Antes de iniciar, aseg\'urese de tener instalado:

\begin{itemize}[leftmargin=2em]
    \item \textbf{Python 3.8+}: Verifique con \texttt{python --version}
    \item \textbf{Flask}: Solo para la interfaz web. Instale con \texttt{pip install flask}
\end{itemize}

\subsection{Inicio R\'apido --- Interfaz Web (Recomendada)}

\begin{enumerate}[leftmargin=2em]
    \item Abra una terminal (PowerShell o CMD).
    \item Navegue al directorio del proyecto:
    \begin{lstlisting}
cd "proyecto final paradigmas"
    \end{lstlisting}
    \item Ejecute el servidor web:
    \begin{lstlisting}
python app.py
    \end{lstlisting}
    \item Ver\'a el siguiente mensaje:
    \begin{lstlisting}[style=console]
  === SISTEMA DE GESTION DE MANTENIMIENTO ===
  Interfaz Web -> http://localhost:5000
  ==========================================

 * Running on http://127.0.0.1:5000
    \end{lstlisting}
    \item Abra su navegador y vaya a \texttt{http://localhost:5000}
    \item Para detener el servidor, presione \texttt{Ctrl+C} en la terminal.
\end{enumerate}

\subsection{Inicio --- Interfaz de Consola}

\begin{lstlisting}
python main.py
\end{lstlisting}

Se mostrar\'a el men\'u principal directamente en la terminal.

\begin{warnbox}[Primera ejecuci\'on]
La primera vez que ejecute el sistema, se crear\'a autom\'aticamente el archivo \texttt{mantenimiento.db} con la base de datos y los datos de ejemplo precargados (15 equipos, 12 t\'ecnicos, 15 \'ordenes de trabajo).
\end{warnbox}

% ══════════════════════════════════════════════════════════════
% SECCIÓN 2: MENÚ PRINCIPAL — INTERFAZ WEB
% ══════════════════════════════════════════════════════════════
\newpage
\section{Interfaz Web --- Navegaci\'on}

La interfaz web cuenta con un \textbf{panel lateral (sidebar)} con las siguientes secciones:

\begin{table}[H]
    \centering
    \begin{tabularx}{\textwidth}{l l X}
        \toprule
        \textbf{Icono} & \textbf{Secci\'on} & \textbf{Descripci\'on} \\
        \midrule
        Gr\'afico & \textbf{Dashboard} & P\'agina principal con estad\'isticas generales, \'ordenes atrasadas y \'ultimas OTs. \\
        \addlinespace
        Llave & \textbf{Equipos} & Listado completo de equipos. Crear, editar y eliminar equipos. \\
        \addlinespace
        Persona & \textbf{T\'ecnicos} & Listado de t\'ecnicos. Crear, editar y eliminar t\'ecnicos. \\
        \addlinespace
        Lista & \textbf{\'Ordenes de Trabajo} & Listado de OTs. Crear, editar, completar, cancelar y eliminar \'ordenes. \\
        \addlinespace
        Documento & \textbf{Historial} & Historial de OTs completadas con filtros por equipo, t\'ecnico y fechas. \\
        \addlinespace
        Barras & \textbf{Reportes} & 4 reportes especiales con consultas avanzadas. \\
        \bottomrule
    \end{tabularx}
    \caption{Secciones del men\'u lateral}
\end{table}

\subsection{Dashboard}

El Dashboard muestra un resumen general del sistema:

\begin{itemize}[leftmargin=2em]
    \item \textbf{Tarjetas de estad\'isticas}: Total de equipos, t\'ecnicos, OTs completadas, OTs activas y OTs atrasadas.
    \item \textbf{Tabla de \'ordenes atrasadas}: OTs que llevan m\'as de 7 d\'ias sin completarse. Incluye d\'ias transcurridos y prioridad.
    \item \textbf{\'Ultimas 5 \'ordenes}: Las OTs m\'as recientes con su estado actual.
\end{itemize}

% ══════════════════════════════════════════════════════════════
% SECCIÓN 3: MENÚ PRINCIPAL — CONSOLA
% ══════════════════════════════════════════════════════════════
\newpage
\section{Interfaz de Consola --- Opciones del Men\'u}

Al iniciar con \texttt{python main.py}, se muestra:

\begin{lstlisting}[style=console]
  ============================================
     SISTEMA DE GESTION DE MANTENIMIENTO
  ============================================
  1. Gestion de Equipos
  2. Gestion de Tecnicos
  3. Gestion de Ordenes de Trabajo
  4. Historial de Mantenimiento
  5. Reportes y Consultas Especiales
  0. Salir

  Seleccione una opcion:
\end{lstlisting}

\subsection{Opci\'on 1: Gesti\'on de Equipos}

\begin{table}[H]
    \centering
    \begin{tabularx}{\textwidth}{c l X}
        \toprule
        \textbf{Opci\'on} & \textbf{Acci\'on} & \textbf{Descripci\'on} \\
        \midrule
        1 & Listar todos & Muestra una tabla con todos los equipos registrados. \\
        2 & Buscar por c\'odigo & Busca y muestra el detalle de un equipo espec\'ifico. \\
        3 & Crear nuevo & Solicita los datos para registrar un nuevo equipo. \\
        4 & Actualizar & Permite modificar campos de un equipo existente. \\
        5 & Eliminar & Elimina un equipo (solo si no tiene OTs asociadas). \\
        0 & Volver & Regresa al men\'u principal. \\
        \bottomrule
    \end{tabularx}
    \caption{Submen\'u de Gesti\'on de Equipos}
\end{table}

\subsection{Opci\'on 2: Gesti\'on de T\'ecnicos}

\begin{table}[H]
    \centering
    \begin{tabularx}{\textwidth}{c l X}
        \toprule
        \textbf{Opci\'on} & \textbf{Acci\'on} & \textbf{Descripci\'on} \\
        \midrule
        1 & Listar todos & Muestra una tabla con todos los t\'ecnicos. \\
        2 & Buscar por documento & Busca un t\'ecnico por su n\'umero de documento. \\
        3 & Crear nuevo & Solicita los datos para registrar un t\'ecnico. \\
        4 & Actualizar & Permite modificar datos de un t\'ecnico existente. \\
        5 & Eliminar & Elimina un t\'ecnico (solo si no tiene OTs activas). \\
        0 & Volver & Regresa al men\'u principal. \\
        \bottomrule
    \end{tabularx}
    \caption{Submen\'u de Gesti\'on de T\'ecnicos}
\end{table}

\subsection{Opci\'on 3: Gesti\'on de \'Ordenes de Trabajo}

\begin{table}[H]
    \centering
    \begin{tabularx}{\textwidth}{c l X}
        \toprule
        \textbf{Opci\'on} & \textbf{Acci\'on} & \textbf{Descripci\'on} \\
        \midrule
        1 & Listar todas & Muestra todas las \'ordenes con estado y prioridad. \\
        2 & Ver detalle & Muestra informaci\'on completa de una OT espec\'ifica. \\
        3 & Crear nueva & Crea una nueva orden asignando equipo y t\'ecnico. \\
        4 & Actualizar & Modifica descripci\'on, prioridad o t\'ecnico. \\
        5 & Completar OT & Registra soluci\'on, costo y cierra la orden. \\
        6 & Cancelar OT & Cancela una orden pendiente o en proceso. \\
        7 & Cambiar estado & Cambia manualmente el estado de una OT. \\
        8 & Eliminar & Elimina una OT (solo pendientes o canceladas). \\
        0 & Volver & Regresa al men\'u principal. \\
        \bottomrule
    \end{tabularx}
    \caption{Submen\'u de Gesti\'on de \'Ordenes de Trabajo}
\end{table}

\subsection{Opci\'on 4: Historial de Mantenimiento}

Muestra las \'ordenes de trabajo \textbf{completadas} con opciones de filtrado:

\begin{table}[H]
    \centering
    \begin{tabularx}{\textwidth}{c l X}
        \toprule
        \textbf{Opci\'on} & \textbf{Filtro} & \textbf{Descripci\'on} \\
        \midrule
        1 & Ver todo & Muestra todo el historial sin filtros. \\
        2 & Por equipo & Filtra por c\'odigo de equipo (ej: EQ-001). \\
        3 & Por t\'ecnico & Filtra por documento del t\'ecnico. \\
        4 & Por rango de fechas & Filtra entre fecha inicio y fecha fin. \\
        0 & Volver & Regresa al men\'u principal. \\
        \bottomrule
    \end{tabularx}
    \caption{Submen\'u de Historial}
\end{table}

\subsection{Opci\'on 5: Reportes y Consultas Especiales}

\begin{table}[H]
    \centering
    \begin{tabularx}{\textwidth}{c l X}
        \toprule
        \textbf{Opci\'on} & \textbf{Reporte} & \textbf{Descripci\'on} \\
        \midrule
        1 & Equipos sin mantenimiento & Equipos que nunca han tenido una OT completada. \\
        2 & Top equipos intervenidos & Los 2 equipos con m\'as intervenciones. \\
        3 & T\'ecnico con mayor costo & T\'ecnico con mayor costo acumulado en repuestos. \\
        4 & \'Ordenes atrasadas & OTs pendientes o en proceso con m\'as de 7 d\'ias. \\
        0 & Volver & Regresa al men\'u principal. \\
        \bottomrule
    \end{tabularx}
    \caption{Submen\'u de Reportes}
\end{table}

% ══════════════════════════════════════════════════════════════
% SECCIÓN 4: EJEMPLOS DE USO
% ══════════════════════════════════════════════════════════════
\newpage
\section{Ejemplos de Uso}

\subsection{Ejemplo 1: Crear un Nuevo Equipo}

\subsubsection{En la interfaz web:}
\begin{enumerate}[leftmargin=2em]
    \item Haga clic en \textbf{Equipos} en el men\'u lateral.
    \item Haga clic en el bot\'on \textbf{``+ Nuevo Equipo''}.
    \item Complete el formulario:
    \begin{itemize}
        \item C\'odigo: \texttt{EQ-016}
        \item Nombre: \texttt{Taladro Industrial Bosch}
        \item Tipo: \texttt{Taladro}
        \item Ubicaci\'on: \texttt{Taller Principal - Zona 3}
        \item Fecha de instalaci\'on: \texttt{2026-01-15}
        \item Estado: \texttt{Operativo}
    \end{itemize}
    \item Haga clic en \textbf{``Crear Equipo''}.
    \item Aparecer\'a un mensaje verde de confirmaci\'on.
\end{enumerate}

\subsubsection{En la consola:}
\begin{lstlisting}[style=console]
  Seleccione una opcion: 1      (Gestion de Equipos)
  Seleccione una opcion: 3      (Crear nuevo)

  Codigo del equipo: EQ-016
  Nombre: Taladro Industrial Bosch
  Tipo: Taladro
  Ubicacion: Taller Principal - Zona 3
  Fecha de instalacion (YYYY-MM-DD): 2026-01-15
  Estado (operativo/en mantenimiento/fuera de servicio): operativo

  [OK] Equipo 'EQ-016' creado exitosamente.
\end{lstlisting}

\subsection{Ejemplo 2: Crear una Orden de Trabajo}

\subsubsection{En la interfaz web:}
\begin{enumerate}[leftmargin=2em]
    \item Haga clic en \textbf{\'Ordenes de Trabajo} en el men\'u lateral.
    \item Haga clic en \textbf{``+ Nueva Orden''}.
    \item Complete:
    \begin{itemize}
        \item Fecha de solicitud: \texttt{2026-05-24} (se llena autom\'aticamente).
        \item Equipo: Seleccione de la lista desplegable.
        \item T\'ecnico: Seleccione de la lista desplegable.
        \item Prioridad: \texttt{Alta}, \texttt{Media} o \texttt{Baja}.
        \item Descripci\'on de la falla: Describa el problema detectado.
    \end{itemize}
    \item Haga clic en \textbf{``Crear Orden''}.
\end{enumerate}

\subsubsection{En la consola:}
\begin{lstlisting}[style=console]
  Seleccione una opcion: 3      (Gestion de OTs)
  Seleccione una opcion: 3      (Crear nueva)

  Fecha de solicitud (YYYY-MM-DD) [Enter=hoy]: 
  Codigo del equipo: EQ-016
  Documento del tecnico: 1001234567
  Prioridad (alta/media/baja): media
  Descripcion de la falla: Vibracion excesiva en el husillo

  [OK] Orden de trabajo #16 creada exitosamente.
\end{lstlisting}

\subsection{Ejemplo 3: Completar una Orden de Trabajo}

Este es uno de los flujos m\'as importantes del sistema, ya que aplica la regla de negocio de cambiar el equipo a estado ``operativo''.

\subsubsection{En la interfaz web:}
\begin{enumerate}[leftmargin=2em]
    \item En la lista de \'ordenes, haga clic en \textbf{``Completar''} junto a la OT deseada.
    \item Ingrese la \textbf{soluci\'on aplicada}: descripci\'on de lo que se hizo.
    \item Ingrese el \textbf{costo de repuestos}: valor num\'erico.
    \item Haga clic en \textbf{``Completar Orden''}.
    \item El sistema autom\'aticamente:
    \begin{itemize}
        \item Registra la fecha de ejecuci\'on (hoy).
        \item Cambia el estado de la OT a ``completada''.
        \item Cambia el equipo asociado a estado ``operativo''.
    \end{itemize}
\end{enumerate}

\subsubsection{En la consola:}
\begin{lstlisting}[style=console]
  Seleccione una opcion: 3      (Gestion de OTs)
  Seleccione una opcion: 5      (Completar OT)

  ID de la orden a completar: 9
  Solucion aplicada: Reemplazo de impulsor y alineacion
  Costo de repuestos: 350000

  [OK] Orden #9 completada. Equipo 'EQ-003' ahora esta 
  en estado 'operativo'.
\end{lstlisting}

\subsection{Ejemplo 4: Consultar Reportes Especiales}

\subsubsection{Reporte: Equipos sin Mantenimiento}

\textbf{En la web:} Haga clic en \textbf{Reportes} en el men\'u lateral. El primer reporte muestra los equipos que nunca han tenido una OT completada.

\textbf{En la consola:}
\begin{lstlisting}[style=console]
  Seleccione una opcion: 5      (Reportes)
  Seleccione una opcion: 1      (Equipos sin mantenimiento)

  Equipos que nunca recibieron mantenimiento: 8
  -----------------------------------------------
  EQ-003  Bomba Centrifuga Grundfos
  EQ-008  Prensa Hidraulica Enerpac
  EQ-010  Sistema HVAC Carrier
  EQ-011  Fresadora CNC DMG Mori
  ...
\end{lstlisting}

\subsubsection{Reporte: T\'ecnico con Mayor Costo}

\textbf{En la web:} Se muestra como una tarjeta destacada con el nombre del t\'ecnico, su costo total acumulado y el n\'umero de \'ordenes completadas.

\textbf{En la consola:}
\begin{lstlisting}[style=console]
  Seleccione una opcion: 5      (Reportes)
  Seleccione una opcion: 3      (Tecnico mayor costo)

  Tecnico con mayor costo acumulado:
  ----------------------------------
  Nombre:     Diego Alejandro Herrera
  Documento:  1004445566
  Costo total: $2,500,000.00
  Ordenes completadas: 1
\end{lstlisting}

\subsection{Ejemplo 5: Filtrar el Historial}

\subsubsection{En la interfaz web:}
\begin{enumerate}[leftmargin=2em]
    \item Haga clic en \textbf{Historial} en el men\'u lateral.
    \item Use la barra de filtros:
    \begin{itemize}
        \item Seleccione un \textbf{equipo} espec\'ifico (o ``Todos'').
        \item Seleccione un \textbf{t\'ecnico} espec\'ifico (o ``Todos'').
        \item Ingrese \textbf{fecha desde} y \textbf{fecha hasta}.
    \end{itemize}
    \item Haga clic en \textbf{``Filtrar''}.
    \item Para limpiar filtros, haga clic en \textbf{``Limpiar''}.
\end{enumerate}

\subsubsection{En la consola:}
\begin{lstlisting}[style=console]
  Seleccione una opcion: 4      (Historial)
  Seleccione una opcion: 4      (Filtrar por rango de fechas)

  Fecha inicio (YYYY-MM-DD): 2025-01-01
  Fecha fin (YYYY-MM-DD): 2025-06-30

  Historial de mantenimiento (6 registros):
  ------------------------------------------
  #1  2025-01-10  EQ-001  Compresor Industrial Atlas
  #2  2025-02-05  EQ-002  Torno CNC Haas ST-10
  ...
\end{lstlisting}

% ══════════════════════════════════════════════════════════════
% SECCIÓN 5: CASOS DE USO COMUNES
% ══════════════════════════════════════════════════════════════
\newpage
\section{Casos de Uso Comunes}

\subsection{Caso 1: Registrar una Falla y Asignar un T\'ecnico}

\begin{enumerate}[leftmargin=2em]
    \item Un operario detecta una falla en un equipo.
    \item El supervisor ingresa al sistema y crea una \textbf{nueva orden de trabajo}:
    \begin{itemize}
        \item Selecciona el equipo afectado.
        \item Selecciona el t\'ecnico especializado disponible.
        \item Describe la falla y asigna la prioridad.
    \end{itemize}
    \item El equipo cambia autom\'aticamente a estado ``en mantenimiento''.
    \item La OT aparece en el Dashboard como ``pendiente''.
\end{enumerate}

\subsection{Caso 2: Completar un Mantenimiento}

\begin{enumerate}[leftmargin=2em]
    \item El t\'ecnico termina la reparaci\'on.
    \item El supervisor busca la OT en el sistema y hace clic en \textbf{``Completar''}.
    \item Ingresa la soluci\'on aplicada y el costo de repuestos.
    \item El sistema:
    \begin{itemize}
        \item Registra la fecha de ejecuci\'on.
        \item Cambia la OT a ``completada''.
        \item Cambia el equipo a ``operativo'' autom\'aticamente.
    \end{itemize}
    \item La OT ahora aparece en el \textbf{Historial de Mantenimiento}.
\end{enumerate}

\subsection{Caso 3: Revisar \'Ordenes Atrasadas}

\begin{enumerate}[leftmargin=2em]
    \item El jefe de mantenimiento abre el \textbf{Dashboard}.
    \item Revisa la tarjeta \textbf{``OTs Atrasadas''} (m\'as de 7 d\'ias).
    \item Identifica las \'ordenes cr\'iticas por prioridad y d\'ias transcurridos.
    \item Hace clic en \textbf{``Ver''} para abrir el detalle de cada una.
    \item Toma acciones correctivas: reasignar t\'ecnico, cambiar prioridad, etc.
\end{enumerate}

\subsection{Caso 4: Generar un Informe de Costos}

\begin{enumerate}[leftmargin=2em]
    \item El administrador abre la secci\'on de \textbf{Reportes}.
    \item Consulta \textbf{``T\'ecnico con Mayor Costo Acumulado''} para identificar al t\'ecnico con mayor gasto en repuestos.
    \item Consulta el \textbf{Historial} filtrando por rango de fechas del mes.
    \item Revisa el costo de repuestos de cada OT completada.
\end{enumerate}

\subsection{Caso 5: Dar de Baja un Equipo}

\begin{enumerate}[leftmargin=2em]
    \item Aseg\'urese de que el equipo \textbf{no tenga \'ordenes de trabajo} asociadas.
    \item Si tiene OTs, primero debe completar o cancelar todas las OTs.
    \item Luego, vaya a \textbf{Equipos} y haga clic en \textbf{``Eliminar''}.
    \item Confirme la eliminaci\'on.
\end{enumerate}

\begin{warnbox}[Regla de negocio]
No es posible eliminar un equipo que tenga \'ordenes de trabajo asociadas. El sistema mostrar\'a un mensaje de error indicando cu\'antas OTs tiene asignadas.
\end{warnbox}

% ══════════════════════════════════════════════════════════════
% SECCIÓN 6: MENSAJES DE ERROR
% ══════════════════════════════════════════════════════════════
\newpage
\section{Mensajes de Error y C\'omo Resolverlos}

\subsection{Errores de Reglas de Negocio}

\begin{errorbox}[Error: No se puede eliminar equipo con OTs]
\textbf{Mensaje:} \textit{``No se puede eliminar: El equipo `EQ-001' tiene 2 orden(es) de trabajo asociada(s).''}

\textbf{Causa:} Se intent\'o eliminar un equipo que tiene \'ordenes de trabajo.

\textbf{Soluci\'on:}
\begin{enumerate}
    \item Complete o cancele todas las \'ordenes asociadas al equipo.
    \item Luego elimine las \'ordenes (solo las canceladas o pendientes).
    \item Finalmente, intente eliminar el equipo nuevamente.
\end{enumerate}
\end{errorbox}

\vspace{0.5cm}

\begin{errorbox}[Error: T\'ecnico no registrado]
\textbf{Mensaje:} \textit{``No se puede asignar: El t\'ecnico con documento `9999999999' no est\'a registrado.''}

\textbf{Causa:} Se intent\'o crear una OT asignando un documento de t\'ecnico que no existe en la base de datos.

\textbf{Soluci\'on:}
\begin{enumerate}
    \item Verifique que el n\'umero de documento sea correcto.
    \item Si el t\'ecnico es nuevo, primero reg\'istrelo en \textbf{Gesti\'on de T\'ecnicos}.
    \item Luego cree la orden de trabajo con el documento correcto.
\end{enumerate}
\end{errorbox}

\vspace{0.5cm}

\begin{errorbox}[Error: Email duplicado]
\textbf{Mensaje:} \textit{``Error al crear el t\'ecnico. Verifique que el documento y email sean \'unicos.''}

\textbf{Causa:} Se intent\'o registrar un t\'ecnico con un email que ya est\'a en uso.

\textbf{Soluci\'on:} Use un email diferente para el nuevo t\'ecnico.
\end{errorbox}

\vspace{0.5cm}

\begin{errorbox}[Error: Estado inv\'alido]
\textbf{Mensaje:} \textit{``Error: Estado inv\'alido. Debe ser: operativo, en mantenimiento, fuera de servicio.''}

\textbf{Causa:} Se ingres\'o un estado que no est\'a en la lista permitida.

\textbf{Soluci\'on:} Use exactamente uno de estos valores: \texttt{operativo}, \texttt{en mantenimiento}, \texttt{fuera de servicio}.
\end{errorbox}

\subsection{Errores de Ejecuci\'on}

\begin{errorbox}[Error: ModuleNotFoundError: flask]
\textbf{Mensaje:} \textit{``ModuleNotFoundError: No module named `flask'''}

\textbf{Causa:} Flask no est\'a instalado.

\textbf{Soluci\'on:}
\begin{lstlisting}
pip install flask
\end{lstlisting}
\end{errorbox}

\vspace{0.5cm}

\begin{errorbox}[Error: UnicodeEncodeError (consola)]
\textbf{Mensaje:} \textit{``UnicodeEncodeError: `charmap' codec can't encode character...''}

\textbf{Causa:} La consola de Windows no soporta caracteres especiales Unicode.

\textbf{Soluci\'on:}
\begin{lstlisting}
python -X utf8 main.py
\end{lstlisting}
\end{errorbox}

\vspace{0.5cm}

\begin{errorbox}[Error: Address already in use (puerto 5000)]
\textbf{Mensaje:} \textit{``OSError: [Errno 98] Address already in use''}

\textbf{Causa:} Otra instancia del servidor ya est\'a corriendo en el puerto 5000.

\textbf{Soluci\'on:}
\begin{enumerate}
    \item Cierre la otra instancia con \texttt{Ctrl+C}.
    \item O busque y termine el proceso: \texttt{taskkill /F /IM python.exe} (Windows).
    \item Vuelva a ejecutar \texttt{python app.py}.
\end{enumerate}
\end{errorbox}

\vspace{0.5cm}

\begin{errorbox}[Error: sqlite3.ProgrammingError (hilos)]
\textbf{Mensaje:} \textit{``sqlite3.ProgrammingError: Los objetos SQLite creados en un hilo solo pueden usarse en ese mismo hilo.''}

\textbf{Causa:} Versi\'on anterior de \texttt{database.py} sin soporte multi-hilo.

\textbf{Soluci\'on:} Aseg\'urese de que \texttt{database.py} use \texttt{check\_same\_thread=False}:
\begin{lstlisting}[language=Python]
self._conexion = sqlite3.connect(
    DB_PATH, check_same_thread=False
)
\end{lstlisting}
\end{errorbox}

% ══════════════════════════════════════════════════════════════
% SECCIÓN 7: REFERENCIA RÁPIDA
% ══════════════════════════════════════════════════════════════
\newpage
\section{Referencia R\'apida}

\subsection{Comandos de Ejecuci\'on}

\begin{table}[H]
    \centering
    \begin{tabularx}{\textwidth}{l X}
        \toprule
        \textbf{Comando} & \textbf{Descripci\'on} \\
        \midrule
        \texttt{python app.py} & Inicia la interfaz web en \texttt{http://localhost:5000} \\
        \texttt{python main.py} & Inicia la interfaz de consola \\
        \texttt{pip install flask} & Instala la dependencia Flask \\
        \texttt{python -X utf8 main.py} & Ejecuta la consola con soporte UTF-8 \\
        \bottomrule
    \end{tabularx}
    \caption{Comandos principales}
\end{table}

\subsection{Estados V\'alidos}

\begin{table}[H]
    \centering
    \begin{tabularx}{\textwidth}{l l X}
        \toprule
        \textbf{Entidad} & \textbf{Estados} & \textbf{Significado} \\
        \midrule
        \multirow{3}{*}{Equipos} & \texttt{operativo} & Equipo funcionando normalmente \\
        & \texttt{en mantenimiento} & Equipo en reparaci\'on o revisi\'on \\
        & \texttt{fuera de servicio} & Equipo deshabilitado \\
        \addlinespace
        \multirow{4}{*}{\'Ordenes} & \texttt{pendiente} & OT creada, sin iniciar \\
        & \texttt{en proceso} & T\'ecnico trabajando en la OT \\
        & \texttt{completada} & OT resuelta exitosamente \\
        & \texttt{cancelada} & OT cancelada antes de completarse \\
        \addlinespace
        \multirow{3}{*}{Prioridades} & \texttt{alta} & Atenci\'on urgente requerida \\
        & \texttt{media} & Atenci\'on programada \\
        & \texttt{baja} & Puede esperar \\
        \bottomrule
    \end{tabularx}
    \caption{Estados y prioridades v\'alidas del sistema}
\end{table}

\subsection{Reglas de Negocio -- Resumen}

\begin{tipbox}[Reglas importantes a recordar]
\begin{enumerate}
    \item \textbf{No eliminar equipos con OTs:} Primero complete o cancele las \'ordenes.
    \item \textbf{T\'ecnico debe existir:} Registre al t\'ecnico antes de asignarlo a una OT.
    \item \textbf{Completar OT = equipo operativo:} Al completar una orden, el equipo asociado cambia autom\'aticamente a estado ``operativo''.
    \item \textbf{Datos no se duplican:} Al re-ejecutar el programa, no se crean datos duplicados.
\end{enumerate}
\end{tipbox}

\end{document}
